\section{Discussion and Limitations}\label{sec:discussion}

\noindent\textbf{Where Should Coordination Matter?} The null E-versus-F result does not contradict the design; it bounds where the design can pay off. An outcome-level benefit requires a regime in which the choice of $\alpha$ changes \emph{which} escape manoeuvres are admitted at outcome-deciding moments: faster or persistent movers, tighter margins, or a dynamic gap the escape route must thread while the barrier is active. Our generator's pedestrian-like movers produce barrier engagement (\cref{tab:slack}) but not that sensitivity. Constructing and measuring such a regime is future work, and we claim no coordination benefit we have not measured.

\noindent\textbf{Cost Analysis.} We do not report per-cycle compute. The benchmark runs a Python mirror whose timings say nothing about the C++ plugin, and the compute telemetry of the committed live pilot trials is empty ($n=0$ samples), so any figure would be unsupported. The QP adds one small solve per cycle whose size grows with the obstacle budget $N_{\mathrm{obs}}$ (capped by clearance pruning), and the escape terms are evaluated only while entrapped; measuring their cost on a real-time-capable host belongs to the full-scale benchmark.

\noindent\textbf{Deployment Lessons.} The three live faults of \cref{sec:live} share a pattern: each sits on an interface (costmap semantics, the compute budget, the binary interface across the pluginlib boundary) that unit tests with one consistent build configuration and synthetic inputs do not exercise. For controller plugins that write into another library's tensors, we recommend compiling against the installed headers' exact flags and committing a regression test that replays live-captured shapes.

\noindent\textbf{Limitations.} We state the limits of the evidence explicitly.
\begin{itemize}
    \item \emph{Full-scale physics benchmark not run.} The quantitative comparison is the 2D mirror, not the live controller. The specified full-scale protocol, an identical Nav2 stack and differential robot across BARN \citep{barn,barn_challenge}, DynaBARN \citep{dynabarn} and HuNavSim \citep{hunavsim} against stock MPPI \citep{nav2,mppi_orig}, DWB \citep{dwa}, RPP \citep{rpp}, TEB \citep{teb}, a DRPA-style always-on escape \citep{drpa_mppi} and a Shield-style CBF-only controller \citep{shield_mppi}, scored by success, collision, time-to-goal, path length, clearance and compute under the same statistics, needs a GPU workstation we did not have. The live entrapment-and-escape reproduction is deferred for the same reason.
    \item \emph{Coordination benefit unmeasured.} The coordination is certified-safe by construction but produced no outcome difference over independent escape+CBF (\cref{sec:ablation}).
    \item \emph{Braking is not a safe fallback when cornered.} In the narrow-dynamic pocket the infeasible-QP fallback ($v=0$) leaves the robot in the path of a sweeping mover, and the CBF configurations collide more often than their no-CBF counterparts ($18\%\to32\%$). A cornered robot needs an evasive fallback, not a stop.
    \item \emph{Look-ahead point versus full body.} \Cref{prop:invariance} certifies the look-ahead point $\mathbf{P}_t$, not the robot body. Body safety is approximate: $L\approx r$ places $\mathbf{P}_t$ at the leading edge, and the costmap critic and \code{collision_monitor} cover the residual. Full-body certification would inflate the margin by $L$, which we avoid because it over-constrains sub-metre gaps. The gap is visible in the data: body-grazing crossers still collide at $16$--$18\%$ with the filter.
    \item \emph{Sampled-data and slack caveats.} Invariance between samples is assumed (\cref{as:intersample}), not proven, and the shared slack couples the rows of one solve; we report slack usage instead of an unconditional certificate.
    \item \emph{Constant-velocity tracker.} The obstacle model evaluated here is a constant-velocity baseline that mispredicts turning or accelerating agents and does not quantify its uncertainty. The released controller has since added occupancy-persistence static/dynamic classification, least-squares horizon prediction and an online conformal radius; these belong to follow-up work, and none of them is present in the 2D mirror behind every quantitative result reported here.
    \item \emph{Other approximations.} The TTC is a one-dimensional closing-speed estimate that can induce over-rotation in tight spaces, and localization jumps in symmetric or narrow maps can perturb the progress signal.
\end{itemize}
